\section{Rare matched walks and continuation records}\label{sec:walks}

Fix a finite set $S\subseteq M_p$ containing $o$. Its points will serve
as states, and its labels and metric distances remain fixed throughout
the following argument. The set may be the support obtained from an
operator in \Cref{lem:finite-support}; its points may lie in any of the
graphs or among the anchors. We will select a graph only after $S$ has
been fixed.

Recall the incidence convention from \Cref{sec:blocks}. Each edge
instance has two distinct oriented incidences $e,\bar e$, exchanged
by reversal, including when the edge is a loop. A loop therefore counts
twice toward degree. Write $e:x\to y$ and let $\mathcal D(x)$ be the
set of $\Delta$ outgoing incidences at $x$. A walk is
\emph{nonbacktracking} if it never immediately takes the reverse
incidence of its preceding step.

In the application to \Cref{lem:measure-bounds}, we will compare the
measures at the ends of an incidence $e:x\to y$ using a $1$-Lipschitz
function that agrees with $b_x$ on one subset of $S$ and with $-b_y$ on
another. This section supplies the combinatorial rule for choosing those
subsets from the fixed graph sequence and the labels and distances on $S$.

A graph vertex $x$ and a state $j\in S$ are \emph{compatible} when
\[
 \norm{h(x)-h(j)}_\infty\le K+2L+1=D-4.
\]
A nonbacktracking walk $x_0,e_1,x_1,\ldots,e_t,x_t$ in one graph
$G_m$ is \emph{matched} if there are states $j_0,\ldots,j_t\in S$
such that $(x_s,j_s)$ is compatible for every $s$ and
\[
 d_p(j_{s-1},j_s)\le2L\qquad(1\le s\le t).
\]
Both the states and the source vertices may repeat. For example, a
return along a different parallel edge is allowed; only the reverse
of the preceding incidence is forbidden.

The compatibility threshold allows for a label distance below $K$ where
a bump is positive, one source-edge step of length at most one, and one
state step of length at most $2L$. The inequality $K+2L+1<D$ will also
restrict the possible source colors.

\subsection{A long match is rare}

Each state has at most $Q$ possible successors, by the radius-$2L$
ball bound in \Cref{lem:block-geometry}, while a graph walk encounters
many possible colors. We use this disparity to bound the
fraction of roots from which a long matched walk starts.

\begin{lemma}\label{lem:rare-walks}
For a fixed integer $n\ge1$, let $\beta_m$ be the fraction of vertices
of $G_m$ from which a matched walk of length $n$ starts. Then
\begin{equation}\label{eq:bad-fraction}
 \limsup_{m\to\infty}\beta_m
 \le |S|\frac Hk\left(\frac{\Delta QH}{k}\right)^n.
\end{equation}
Consequently one can choose $n\ge2$ and then $m$ so that
$\beta_m<1/(20V)$.
\end{lemma}

\begin{proof}
For $j\in S$, set
\[
 C_j=\{i:h_i(j)<D\}.
\]
\Cref{lem:block-geometry} gives $|C_j|\le H$. If $x$ has color $i$
and is compatible with $j$, then $h_i(x)=0$ implies
\[
 h_i(j)\le K+2L+1=D-4<D.
\]
Thus the color of $x$ belongs to $C_j$. This is only a necessary
condition for compatibility; it avoids any independence assertion
about the full labels.

Consider the weaker event that a nonbacktracking length-$n$ walk
$x_0,e_1,x_1,\ldots,e_n,x_n$ from the root has states
$j_0,\ldots,j_n\in S$ with $d_p(j_{s-1},j_s)\le2L$ for $1\le s\le n$
and with the color of $x_s$ in $C_{j_s}$ for $0\le s\le n$.
Its occurrence is determined
by the induced colored radius-$n$ ball: every traversed vertex and
edge lies in that ball, including loops and parallel edge instances.
There are only finitely many such rooted types. With $S$ and $n$ fixed,
the sets $C_j$ and the transition relation are fixed, so the event selects
the same finite union of these types for every $m$.
\Cref{lem:local-limit} gives convergence for
every type and therefore for this particular finite union. Its empirical
frequency converges to its probability on the independently colored
$\Delta$-regular tree. The support may have been chosen using the entire
deterministic graph sequence: once fixed, it only selects the types in
the union. No random colors are conditioned on the support.

On the tree there are at most $\Delta^n$ nonbacktracking walks of
length $n$ from the root, each visiting $n+1$ distinct vertices.
For each state $j$, its allowed successors lie in
$S\cap\overline B_{d_p}(j,2L)$, which has at most $Q$ points.
Thus there are at most $|S|Q^n$ allowed state sequences. For a fixed
walk and state sequence, independence of the colors gives probability
at most $(H/k)^{n+1}$ for all required memberships. Repeated states
and overlapping sets $C_j$ do not affect this product, because the
source vertices on the tree are distinct. A union bound gives
\[
 \Delta^n|S|Q^n(H/k)^{n+1}
 =|S|(H/k)(\Delta QH/k)^n.
\]
Every matched walk satisfies the weaker event, proving
\eqref{eq:bad-fraction}. Finally, $k>2\Delta QH$ gives
$\Delta QH/k<1/2$. Choose $n\ge2$ so that the right-hand side of
\eqref{eq:bad-fraction} is less than $1/(40V)$, and then choose $m$
with $\beta_m<1/(20V)$.
\end{proof}

\begin{remark}[Adding roots with non-tree neighborhoods]\label{rem:non-tree-bad-set}
One may enlarge the set counted by $\beta_m$ to include every root whose
induced radius-$(n+1)$ neighborhood is not a tree. For fixed $n$, let
$\gamma_m$ be the fraction of roots with such a neighborhood and
$\widetilde\beta_m$ the enlarged fraction. Then
$\beta_m\le\widetilde\beta_m\le\beta_m+\gamma_m$, while the local limit
gives $\gamma_m\to0$. Thus \eqref{eq:bad-fraction} remains valid with
$\widetilde\beta_m$ in place of $\beta_m$, and the same successive choices
of $n\ge2$ and sufficiently large $m$ give
$\widetilde\beta_m<1/(20V)$.
\end{remark}

\subsection{At most one exceptional direction per state}

Fix $n\ge2$ and one graph $G=G_m$. Call a vertex \emph{bad} if a
matched walk of length $n$ starts there, and \emph{good} otherwise.
The next argument is deterministic. At a good vertex it uses the
absence of long matched walks to select at most one direction for
each state.

A related comparison of nonbacktracking continuation lengths appears in
Ben-David, Childs, Gily\'en, Kretschmer, Podder, and Wang
\cite[\S6.1.2, Algorithm~6.1 and Lemma~6.3]{BenDavidEtAl24}. Their
probabilistic graph algorithm uses such lengths to select parent directions,
and their consistency lemma excludes an incompatible pair of choices by
symmetric length inequalities. Here the
records are deterministic and indexed by support states; the incidence
formulation also includes loops and parallel edges. The argument below is
self-contained.

For a good vertex $x$, an incidence $e\in\mathcal D(x)$, and a
state $j\in S$, let $\mathcal L_x(e,j)$ be the set of positive integer
lengths of matched walks starting through $e$ with initial state $j$.
Define
\[
 r_x(e,j)=\max\bigl(\{0\}\cup\mathcal L_x(e,j)\bigr).
\]
A matched walk of length at least $n$
has a matched prefix of length $n$, so
$\mathcal L_x(e,j)\subseteq\{1,\ldots,n-1\}$. Thus the maximum is
well-defined. The value zero means only that no positive-length walk
exists; it makes no assertion about compatibility of $(x,j)$.

Declare $e$ \emph{exceptional at $(x,j)$} if its record is strictly
greater than every other outgoing record at $(x,j)$. There is at most
one such incidence; if largest records are tied, there is none. At a
bad vertex, declare all outgoing incidences exceptional at every state.
For an incidence $e:x\to y$, set
\[
 E_e=\{j\in S:e\text{ is exceptional at }(x,j)\}.
\]

\begin{lemma}\label{lem:exceptional-directions}
Let $e:x\to y$ be an incidence of $G$ and let $j,j'\in S$ be states.
If $(x,j)$ and $(y,j')$
are compatible and $d_p(j,j')\le2L$, then
$j\in E_e$ or $j'\in E_{\bar e}$.
\end{lemma}

\begin{proof}
If either endpoint is bad, the conclusion follows from the definition.
Suppose both are good, and put
\[
 A_x=\max_{g\in\mathcal D(x)\setminus\{e\}}r_x(g,j),
 \qquad
 A_y=\max_{f\in\mathcal D(y)\setminus\{\bar e\}}r_y(f,j').
\]
The maxima are over nonempty finite sets, since $\Delta\ge4$. If
$A_y>0$, choose a walk
realizing it and prepend $e$, with the state transition $j$ to $j'$.
The old first incidence is not $\bar e$, so the resulting walk is
nonbacktracking and matched. If $A_y=0$, the edge alone, with states
$j,j'$, is a matched walk. The same argument after reversing $e$
therefore gives both inequalities
\begin{equation}\label{eq:record-inequalities}
 r_x(e,j)\ge1+A_y,\qquad r_y(\bar e,j')\ge1+A_x.
\end{equation}
If neither incidence were exceptional, their records would be at
most $A_x$ and $A_y$, respectively. These inequalities would force
\[
 A_x\ge1+A_y\ge2+A_x,
\]
a contradiction.
\end{proof}

The proof uses incidence identities, so it also applies to loops and
parallel edges. It allows repeated states and vertex revisits, and using
the edge by itself covers zero records. The conclusion will let us
prescribe opposite scalar values at the two endpoints: whenever nearby
states could receive opposing nonzero values, at least one of their
prescriptions must be discarded.
